% !TeX root = ../main.tex
% !TeX encoding = utf8

\chapter{Mecanismos de coordinación y diseño de puestos}
% Responsable: Jesús Rodríguez González
% Secciones del índice:
%   3. Principales mecanismos de coordinación
%   4.1 Diseño de puestos (especialización y formalización — inicio)

\section{Principales mecanismos de coordinación}

% Indicar las distintas actividades que se realizan en la sucursal,
% quiénes están implicados en su realización y cómo se coordinan.
%
% Mecanismos a identificar y describir:
%   - Supervisión directa
%   - Normalización de procesos
%   - Normalización de resultados
%   - Normalización de habilidades
%   - Adaptación mutua

\section{Diseño de puestos: especialización}

% 4.1 — Especialización horizontal y vertical
%
% Para cada puesto (o grupo de puestos con características similares):
%   - Especialización horizontal: ¿realiza pocas tareas repetitivas o una gran variedad?
%   - Especialización vertical: ¿tiene autonomía o debe seguir instrucciones estrictas?
%   - Identificar si es un puesto profesional (ampliado verticalmente).

\section{Diseño de puestos: formalización}

% 4.1 — Formalización
%
%   - ¿Están documentadas por escrito las tareas de cada puesto?
%   - ¿Existe un job description formal?
%   - ¿En qué medida las situaciones del día a día están reguladas por normas y protocolos?
%   - ¿Cómo se formaliza el trabajo (manuales, sistemas técnicos, supervisión...)?

\endinput
% -------------------------------------------------------------------
% FIN DEL CAPÍTULO — Jesús Rodríguez González
% -------------------------------------------------------------------
